\chapter{Running the Station day to day}\label{ch:running}

With the program installed, the Station prepared and the config set, day to day is: launch, run a
Test Run, and know where the Test Report went. How the screens look is deliberately not described
here --- the interface is still being worked on --- only what the application does.

\section{Launching}

\begin{lstlisting}[language=ATEShell]
ate
ate --config /tmp/station/config/config.json
ate --allow-unsandboxed-pdf
\end{lstlisting}

With no subcommand, ATE opens the application. At launch, in order, it reads the config; creates the
log and reports directories if they do not exist (refusing with the path if it cannot); opens this
launch's application log; checks the database's schema version and refuses if it is not this ATE's;
closes any Test Run left open by a previous crash; and decides, once, how PDFs will be handled for
this launch.

An invalid config or an unusable database is a dialog with one or two lines and an exit code of 1.
There is no traceback, and nothing is migrated or repaired silently.

\section{Reports, and the PDF sandbox question}

A Test Report is written when a Test Run finishes, under \texttt{reports.directory}, in a
subdirectory named after the UUT Type:

\begin{lstlisting}[language=ATEShell]
/var/lib/ate/reports/RailGun/RailGun_4471_20261101-091203_PASS.html
\end{lstlisting}

The HTML is always written. The PDF is not, and this is the part of ATE that surprises people.
ATE never renders a PDF in its own process: it starts a hidden child, \texttt{ate \_pdf <html>
<pdf> --paper A4}, and waits for it. Chromium's renderer sandbox will not start under root or an
administrator, so a privileged ATE has to be told to give up that isolation, and ATE never does it
by itself:

\begin{center}
\footnotesize
\begin{tabular}{@{}>{\raggedright\arraybackslash}p{4.2cm}>{\raggedright\arraybackslash}p{10.6cm}@{}}
\toprule
Situation & What happens \\
\midrule
Ordinary user & The PDF is rendered with the sandbox on. Nothing is asked. Giving
  \texttt{--allow-unsandboxed-pdf} here does nothing, and ATE logs that it had no effect. \\
Privileged, no flag & A dialog asks, per PDF: "Render PDF without sandbox?", with buttons
  \texttt{Generate PDF without sandbox} and \texttt{Skip PDF}. \texttt{Skip PDF} is the default
  button. Saying no leaves the HTML and writes no PDF. \\
Privileged, \texttt{--allow-}\allowbreak\texttt{unsandboxed-pdf} & No dialog; the PDF is rendered unsandboxed and the
  choice is logged as a WARNING. \\
Privileged, headless & There is nobody to ask, so the PDF is skipped and the log says so --- unless
  \texttt{--allow-unsandboxed-pdf} was given. \\
\bottomrule
\end{tabular}
\end{center}

The flag belongs to \texttt{ate}, so it may come before a hidden subcommand:
\texttt{ate --allow-unsandboxed-pdf \_selftest --run} is the same decision as putting it at the end.
It applies to every Test Report PDF of that launch, including one regenerated from History.

Two things a helper can fail on that are worth knowing:

\begin{itemize}
  \item It refuses to render as root or an administrator without \texttt{--unsandboxed}, because
    that is exactly the choice ATE will not make for you. You only see this if you run
    \texttt{ate \_pdf} yourself.
  \item A PDF can be a perfectly valid PDF with nothing in it. Debian containers ship no fonts; with
    no font, Chromium draws no glyphs and writes about 1.7 kilobytes of empty page. \texttt{ate
    \_selftest --run} checks for text in the PDF precisely so this is caught as a packaging failure
    rather than passed as a success. If a Test Report opens as a blank page, install fontconfig and
    a font before blaming ATE.
\end{itemize}

Environment inherited from another program cannot switch the sandbox off for ATE's PDFs. A
\texttt{QTWEBENGINE\_}\allowbreak\texttt{DISABLE\_SANDBOX}, or a \texttt{--no-sandbox} in a
\texttt{QTWEBENGINE\_}\allowbreak\texttt{CHROMIUM\_FLAGS}, is stripped from the helper's environment
and logged as a WARNING, because nobody chose it. The only way to get an unsandboxed render is the
flag or the dialog.

\section{History, and the Run log}

History is one flat table of Test Runs, newest first, sortable and filterable by UUT Type, serial
substring, result, End Reason, date range and Project. Rows load page by page. It is read-only:
Test Runs are never edited or deleted, so there is no "clean up the database" button, and archiving
means copying the database file.

For a finished run --- including one recorded as \texttt{INTERRUPTED} --- History can open its
report, regenerate the HTML or the PDF with or without the Run log, and export the Run log as CSV,
JSON Lines or text. Exporting is not available for a run that is still in progress. Browsing
History while a run is running is allowed unless \texttt{ui.allowHistoryDuringRun} is set to false.

\cando{run a Test Run, find the Test Report it wrote, and explain why a privileged Station wrote no
PDF --- and what to do about it.}
